\section{Mathematical Framework}

\subsection{Distributional Derivatives}

In the following let $\Omega\subset\R^d,\, d\ge 1$ denote a Lebesgue measurable open set with non-empty interior and let $m\ge 1$.
\begin{definition}[Partial Derivative, Gradient, Jacobian and Divergence]
	Let $f:\Omega\to\R$ and $F:\Omega\to\R^m$ be sufficiently regular.
	\begin{itemize}
		\item[i)] The {\tt partial derivative} of $f$ with respect to $\alpha\in\N_0^d$ is
		      \begin{equation*}
			      \Diff^\alpha f := \frac{\partial^{|\alpha|}f}{\partial x_1^{\alpha_1}\cdots\partial x_d^{\alpha_d}}:\Omega\to\R.
		      \end{equation*}

		\item[ii)] The {\tt gradient} of $u$ is
		      \begin{equation*}
			      \nabla f := \left(\frac{\partial f}{\partial x_1},\ldots,\frac{\partial f}{\partial x_d}\right)^\top:\Omega\to\R^d.
		      \end{equation*}

		\item[iii)] The {\tt Jacobian} (or derivative) of $F=(F_1,\ldots,F_m)^\top$ is
		      \begin{equation*}
			      \Diff F := \left(\frac{\partial F_j}{\partial x_k}\right)_{1\le j\le m}^{1\le k\le d}:\Omega\to\R^{m\times d}.
		      \end{equation*}

		\item[iv)] The {\tt divergence} of $F=(F_1,\ldots,F_m)^\top$ is
		      \begin{equation*}
			      \diver F := \sum_{j=1}^d \frac{\partial F_j}{\partial x_j}:\Omega\to\R.
		      \end{equation*}
	\end{itemize}
\end{definition}

\begin{definition}[Spaces of Continuously Differentiable Functions]
	We denote by $\cC^0(\Omega;\R^m)$ the space of continuous functions $f:\Omega\to\R^m$ and by
	\begin{equation*}
		\cC^k(\Omega;\R^m):= \left\{f\in \cC^0(\Omega;\R^m)\;\middle|\; |\alpha|\le k\;\Rightarrow\;\Diff^\alpha f\in\cC^0(\Omega;\R^m)\right\},\quad k\ge 0,
	\end{equation*}
	the space of such functions that are continuously differentiable up to order $k\ge 0$. Moreover,
	\begin{equation*}
		\cC^\infty(\Omega;\R^m) := \bigcap_{k\ge 0}\cC^k(\Omega;\R^m)
	\end{equation*}
	denotes the space of infinitely continuously differentiable functions.
\end{definition}

\begin{definition}[Test function]
	The space $\cD(\Omega;\R^m)$ contains all infinitely often differentiable functions with compact support in $\Omega$, i.e.\
	\begin{equation*}
		\cD(\Omega;\R^m) = \left\{f\in\cC^\infty(\Omega;\R^m)\;\middle|\;\supp(f)\subset\Omega\text{\em\enspace compact}\right\}
	\end{equation*}
	For $m=1$, we abbreviate $\cD(\Omega):= \cD(\Omega;\R^1)$ and note that
	\begin{equation*}
		\cD(\Omega;\R^m) = (\cD(\Omega))^m,\quad \forall\,m\ge 1.
	\end{equation*}
	A function $f\in\cC^\infty_c(\Omega;\R^m)$ is called a {\tt test function}.
\end{definition}

\begin{definition}[Distribition]
	A {\tt distribution} $T$ is a linear continuous functional on the test function space $\cD(\Omega)$, i.e. $T \in\cD^\prime(\Omega)$. For $m\ge 1$ we set
	\begin{equation*}
		\cD^\prime(\Omega;\R^m) := (\cD^\prime(\Omega))^m
	\end{equation*}
	and for $T=(T_1,\ldots,T_m)\in\cD'(\Omega;\R^m)$ and $\varphi=(\varphi_1,\ldots,\varphi_m)\in\cD(\Omega;\R^m)$ we use the duality pairing
	\begin{equation*}
		\langle T,\,\varphi\rangle := \sum_{j=1}^m \langle T_j,\,\varphi_j\rangle.
	\end{equation*}
\end{definition}

\begin{definition}[Partial derivatives of distributions]\label{def:distrpartialderivative}
	Let $T\in\cD'(\Omega)$. For $\alpha\in\N_0^d$ we define the {\tt distributional partial derivative} $\Diff^\alpha T\in\cD'(\Omega)$ of $T$ by
	\begin{equation*}
		\langle \Diff^\alpha T,\,\varphi\rangle := (-1)^{|\alpha|}\,\langle T,\,\Diff^\alpha\varphi\rangle
		\quad\forall\,\varphi\in\cD(\Omega).
	\end{equation*}
	If $T=(T_1,\ldots,T_m)\in\cD'(\Omega;\R^m)$, we set
	\begin{equation*}
		\Diff^\alpha T := (\Diff^\alpha T_1,\ldots,\Diff^\alpha T_m)\in\cD'(\Omega;\R^m).
	\end{equation*}
\end{definition}

\begin{definition}[Locally Integrable Functions]
	The space $L_\mathrm{loc}^1 (\Omega;\R^m)$ of all {\tt locally integrable} functions on $\Omega$ is defined as
	\begin{equation*}
		L_\mathrm{loc}^1 (\Omega;\R^m):= \left\{f:\Omega\to\R^m\text{\em\enspace measurable}\;\middle|\;f|_K\in L^1(K;\R^m)\;\forall K\subset\Omega,\,K\textit{\em\enspace compact}\right\}.
	\end{equation*}
\end{definition}

\begin{definition}[Regular Distribution]
	A {\tt regular distribution} $T_f\in\cD^\prime(\Omega;\R^m)$ is one of the form
	\begin{equation*}
		\langle T_f,\,\varphi\rangle = \int_\Omega f\, \varphi\;\dx,\quad f\in L^1_\mathrm{loc}(\Omega;\R^m).
	\end{equation*}
	In particular, the map
	\begin{equation*}
		\cT: L^1_\mathrm{loc}(\Omega;\R^m)\to\cD^\prime(\Omega;\R^m),\quad f\mapsto T_f
	\end{equation*}
	is injective and $\im\cT\subset\cD^\prime(\Omega;\R^m)$ is the subspace of regular distributions.
\end{definition}

\begin{definition}[Distributional Gradient and Divergence]
	Let $x = (x_1,\ldots, x_d)^\top\in\Omega$.
	\begin{itemize}
		\item[i)] The {\tt distributional gradient} of $T\in\cD^\prime(\Omega)$ is defined by
		      \begin{equation*}
			      \nabla T:= \left(\frac{\partial T}{\partial x_1},\ldots,\frac{\partial T}{\partial x_d}\right)\in\cD^\prime(\Omega;\R^d),
		      \end{equation*}
		      where each partial derivative $\partial T/\partial x_j,\,1\le j\le d$ is taken in the sense of Definition \ref{def:distrpartialderivative} and it holds that
		      \begin{equation*}
			      \langle\nabla T,\,\varphi\rangle = -\langle T,\,\diver\varphi\rangle\quad\forall\,\varphi\in \cD(\Omega;\R^d)
		      \end{equation*}
		\item[ii)] The {\tt distributional divergence} of $T\in\cD^\prime(\Omega;\R^d)$ is defined by
		      \begin{equation*}
			      \diver T:= \sum_{j = 1}^d\frac{\partial T_j}{\partial x_j}\in\cD^\prime(\Omega),\quad T = (T_1,\ldots, T_d).
		      \end{equation*}
		      where each partial derivative $\partial T_j/\partial x_j,\,1\le j\le d$ is taken in the sense of Definition \ref{def:distrpartialderivative} and it holds that
		      \begin{equation*}
			      \langle\diver ,\,\varphi\rangle = -\langle T,\,\nabla \varphi\rangle\quad\forall\,\varphi\in\cD(\Omega).
		      \end{equation*}
	\end{itemize}
\end{definition}

\subsection{Sobolev Spaces}
Throughout, for $u,v:\Omega\to\R^m$ we denote by $u\cdot v$ the Euclidean inner product in $\R^m$ and by $|u|:=\sqrt{u\cdot u}$ the associated Euclidean norm.

\begin{proposition}[$L^2$-Functions are Locally Integrable]
	The space $L^2(\Omega;\R^m)$ of square-integrable functions on $\Omega$ is a subspace of $L^1_\mathrm{loc}(\Omega;\R^m)$. Consequently, every $f\in L^2(\Omega;\R^m)$ is infinitely differentiable in the sense of distributions.
\end{proposition}
\begin{proof}
	Let $K\subset\Omega$ be compact and $f\in L^2(\Omega;\R^m)$. Then $f\in L^2(K;\R^m)$, i.e.\ $\|f\|_{L^2(K;\R^m)}<\infty$ and by Cauchy-Schwarz we have
	\begin{equation*}
		\|f\|_{L^1(K;\R^m)} = \int_K|f|\;\dx\le \left(\int_K|f|^2\;\dx\right)^{1/2}\left(\int_K 1^2\;\dx\right)^{1/2} = \|f\|_{L^2(K;\R^m)}|K|^{1/2}<\infty,
	\end{equation*}
	where $|\cdot|$ denotes the Lebesgue measure on $\R^d$.
\end{proof}

\begin{definition}[Sobolev Spaces of Integer Order]
	Let $\Omega\subset\R^d$ be open, $m\ge 1$, $k\in\N_0$ and $1\le p<\infty$.
	The {\tt Sobolev space} $W^{k,p}(\Omega;\R^m)$ is defined by
	\begin{equation*}
		W^{k,p}(\Omega;\R^m)
		:=\left\{u\in L^p(\Omega;\R^m)\;\middle|\; \Diff^\alpha u\in L^p(\Omega;\R^m)\ \forall\,\alpha\in\N_0^d,\ |\alpha|\le k\right\},
	\end{equation*}
	where $\Diff^\alpha u$ is understood in the sense of distributions. We adopt the usual conventions $H^k(\Omega;\R^m):=W^{k,2}(\Omega;\R^m)$ and $H^0(\Omega;\R^m):= L^2(\Omega;\R^m)$ and denote by
	\begin{equation*}
		(u,\,v)_{0,\Omega}:= (u,\, v)_{L^2(\Omega)} = \int_\Omega u\cdot v\;\dx,\quad u,v\in H^0(\Omega;\R^d).
	\end{equation*}
	the $L^2$-product and analogously the $L^2$-norm by
	\begin{equation*}
		\|u\|_{0,\Omega}:= \|u\|_{L^2(\Omega)} = \left(\int_\Omega |u|^2\;\dx \right)^{1/2},\quad u\in H^0(\Omega;\R^d).
	\end{equation*}
\end{definition}

\begin{proposition}[$H^k(\Omega;\R^m)$ is a Hillbert space]
	The Sobolev space $H^k(\Omega;\R^m)$ together with the inner product
	\begin{equation*}
		(u,\, v)_{k,\Omega}:= \sum_{|\alpha|\le k}(\Diff^\alpha u,\, \Diff^\alpha v)_{0,\Omega},\quad u,v\in H^k(\Omega;\R^m)
	\end{equation*}
	is a Hilbert space and it induces the norm
	\begin{equation*}
		\|u\|_{k,\Omega} = \left(\sum_{|\alpha|\le k}\|\Diff^\alpha u\|_{0,\Omega}^2\right)^{1/2}<\infty,\quad u\in H^k(\Omega;\R^m).
	\end{equation*}\eop
\end{proposition}

\begin{definition}[Sobolev Spaces of Fractional Order]\label{def:fractionalsobolev}
	For $1\le p<\infty$ and $\theta\in(0,1)$ set
	\begin{equation*}
		|u|_{\theta,p,\Omega}:= \left(\int_\Omega\int_\Omega \frac{|u(x)-u(y)|^p}{|x-y|^{d+p\theta}}\;\dx\;\dy\right)^{1/p},\quad u\in W^{k,p}(\Omega;\R^m).
	\end{equation*}
	The {\tt Sobolev space} of order $s = k+\theta \in\R_{>0}\setminus\N,\, k\in \N_0$ is defined by
	\begin{equation*}
		W^{s,p}(\Omega;\R^m):= \left\{u\in W^{k,p}(\Omega;\R^m)\;\middle|\;|\Diff^\alpha u|_{\theta,p,\Omega}<\infty,\;|\alpha|= k\right\}.
	\end{equation*}
	We adopt the usual convention $H^s(\Omega;\R^m):= W^{s,2}(\Omega;\R^m)$.
\end{definition}

\begin{proposition}[$H^s(\Omega;\R^m)$ is a Banach space]
	The Sobolev space $H^s(\Omega;\R^m)$ defined in Definition \ref{def:fractionalsobolev} together with the norm
	\begin{equation*}
		\|u\|_{s,\Omega}:= \left(\|u\|_{k,\Omega}^2+\sum_{|\alpha|=k} |\Diff^\alpha u|^2_{0,2,\Omega}\right)^{1/2},\quad u\in H^s(\Omega;\R^m)
	\end{equation*}
	is a Banach space.\eop
\end{proposition}



\begin{theorem}[Trace Theorem]\label{th:tracetheorem}

\end{theorem}


